% GENERATED FILE. Edit canonical lessons, then run npm run generate:books.
% canonical-source-hash: fnv1a64:08f37cf981016dfc
% canonical-lessons: AR-C92-qitar, AR-C92-jisr, AR-C92-masjid, AR-C92-funduq, AR-C92-maktab

\chapter{Train, Bridge, Mosque, Hotel, Office}
\label{ch:ar-train-bridge-mosque-hotel-office}
\hlchaptermodality{92}

\begin{chapteropening}
\textbf{By the end of this chapter:} \emph{I can say a train, a bridge, a mosque, a hotel and an office.}

Complete the last lesson of chapter 92: I can say a train, a bridge, a mosque, a hotel and an office.
\end{chapteropening}

\section[\texorpdfstring{\ar{قطار} (qiṭār)}{قطار (qiṭār)}]{\ar{قطار} (qiṭār) — a train}

\label{lesson:AR-C92-qitar}



\begin{quote}
Before the new one: say the Arabic for a car, then the Arabic for a bus.
\end{quote}

\subsection*{You'll want to know: \ar{قطار}}
\textbf{\ar{قطار}} — \emph{qiṭār} — \textquotedblleft{}a train\textquotedblright{}.

\begin{cousinweb}
\textbf{\ar{قطار}} (\emph{qiṭār}) is built on the native Arabic root \emph{q-ṭ-r}, \textquotedblleft{}a file of camels\textquotedblright{} (Lane). Lane's \emph{qiṭār} is a string of camels led one behind another; the train is the modern sense.
\end{cousinweb}

\begin{grammarlens}[title={What you've built}]
One more word for this chapter.
\end{grammarlens}

\subsection*{Your turn}
\begin{itemize}
\raggedright
  \item \emph{Say it:} \emph{qiṭār}
  \item \emph{Say it:} \emph{qiṭār}, once more
  \item \emph{Say it:} say \emph{qiṭār}
  \item \emph{Recall:} say the Arabic for a car, then the Arabic for a bus, then say \emph{qiṭār} again
\end{itemize}

\begin{tcolorbox}[breakable,colback=teal!4,colframe=teal!35!black,title={Before you move on}]
What does \ar{قطار} mean? (\textquotedblleft{}A train\textquotedblright{}.) Say it once more.
\end{tcolorbox}
\section[\texorpdfstring{\ar{جسر} (jisr)}{جسر (jisr)}]{\ar{جسر} (jisr) — a bridge}

\label{lesson:AR-C92-jisr}



\begin{quote}
Before the new one: say the Arabic for a bus, then the Arabic for a train.
\end{quote}

\subsection*{You'll want to know: \ar{جسر}}
\textbf{\ar{جسر}} — \emph{jisr} — \textquotedblleft{}a bridge\textquotedblright{}.

\begin{cousinweb}
\textbf{\ar{جسر}} (\emph{jisr}) is built on the native Arabic root \emph{j-s-r}, \textquotedblleft{}a bridge; to be bold\textquotedblright{} (Lane). An old noun; Lane lists it under \emph{j-s-r}.
\end{cousinweb}

\begin{grammarlens}[title={What you've built}]
One more word for this chapter.
\end{grammarlens}

\subsection*{Your turn}
\begin{itemize}
\raggedright
  \item \emph{Say it:} \emph{jisr}
  \item \emph{Say it:} \emph{jisr}, once more
  \item \emph{Say it:} say \emph{jisr}
  \item \emph{Recall:} say the Arabic for a bus, then the Arabic for a train, then say \emph{jisr} again
\end{itemize}

\begin{tcolorbox}[breakable,colback=teal!4,colframe=teal!35!black,title={Before you move on}]
What does \ar{جسر} mean? (\textquotedblleft{}A bridge\textquotedblright{}.) Say it once more.
\end{tcolorbox}
\section[\texorpdfstring{\ar{مسجد} (masjid)}{مسجد (masjid)}]{\ar{مسجد} (masjid) — a mosque}

\label{lesson:AR-C92-masjid}



\begin{quote}
Before the new one: say the Arabic for a train, then the Arabic for a bridge.
\end{quote}

\subsection*{You'll want to know: \ar{مسجد}}
\textbf{\ar{مسجد}} — \emph{masjid} — \textquotedblleft{}a mosque\textquotedblright{}.

\begin{cousinweb}
\textbf{\ar{مسجد}} (\emph{masjid}) is built on the native Arabic root \emph{s-j-d}, \textquotedblleft{}to bow down, to prostrate oneself\textquotedblright{} (Lane). \emph{Masjid} is the \emph{ma-} place noun, \textquotedblleft{}a place of prostration\textquotedblright{}; English \emph{mosque} comes from it by way of Spanish and French.
\end{cousinweb}

\begin{grammarlens}[title={What you've built}]
One more word for this chapter.
\end{grammarlens}

\subsection*{Your turn}
\begin{itemize}
\raggedright
  \item \emph{Say it:} \emph{masjid}
  \item \emph{Say it:} \emph{masjid}, once more
  \item \emph{Say it:} say \emph{masjid}
  \item \emph{Recall:} say the Arabic for a train, then the Arabic for a bridge, then say \emph{masjid} again
\end{itemize}

\begin{tcolorbox}[breakable,colback=teal!4,colframe=teal!35!black,title={Before you move on}]
What does \ar{مسجد} mean? (\textquotedblleft{}A mosque\textquotedblright{}.) Say it once more.
\end{tcolorbox}
\section[\texorpdfstring{\ar{فندق} (funduq)}{فندق (funduq)}]{\ar{فندق} (funduq) — a hotel}

\label{lesson:AR-C92-funduq}



\begin{quote}
Before the new one: say the Arabic for a bridge, then the Arabic for a mosque.
\end{quote}

\subsection*{You'll want to know: \ar{فندق}}
\textbf{\ar{فندق}} — \emph{funduq} — \textquotedblleft{}a hotel\textquotedblright{}.

\begin{cousinweb}
\textbf{\ar{فندق}} (\emph{funduq}) is a loan. Lane reports it as said to be Persian and adds that it seems to come from Greek \emph{pandokheion}, \textquotedblleft{}inn\textquotedblright{}, a word in the Gospel of Luke; EtymDB gives the same Greek source.
\end{cousinweb}

\begin{grammarlens}[title={What you've built}]
One more word for this chapter.
\end{grammarlens}

\subsection*{Your turn}
\begin{itemize}
\raggedright
  \item \emph{Say it:} \emph{funduq}
  \item \emph{Say it:} \emph{funduq}, once more
  \item \emph{Say it:} say \emph{funduq}
  \item \emph{Recall:} say the Arabic for a bridge, then the Arabic for a mosque, then say \emph{funduq} again
\end{itemize}

\begin{tcolorbox}[breakable,colback=teal!4,colframe=teal!35!black,title={Before you move on}]
What does \ar{فندق} mean? (\textquotedblleft{}A hotel\textquotedblright{}.) Say it once more.
\end{tcolorbox}
\section[\texorpdfstring{\ar{مكتب} (maktab)}{مكتب (maktab)}]{\ar{مكتب} (maktab) — an office; a desk}

\label{lesson:AR-C92-maktab}



\begin{quote}
Before the new one: say the Arabic for a mosque, then the Arabic for a hotel.
\end{quote}

\subsection*{You'll want to know: \ar{مكتب}}
\textbf{\ar{مكتب}} — \emph{maktab} — \textquotedblleft{}an office; a desk\textquotedblright{}.

\begin{cousinweb}
\textbf{\ar{مكتب}} (\emph{maktab}) is built on the native Arabic root \emph{k-t-b}, \textquotedblleft{}to write\textquotedblright{} (Lane). \emph{Maktab} is the \emph{ma-} place noun, \textquotedblleft{}a place of writing\textquotedblright{}: a desk, an office.
\end{cousinweb}

\begin{grammarlens}[title={What you've built}]
One more word for this chapter.
\end{grammarlens}

\subsection*{Your turn}
\begin{itemize}
\raggedright
  \item \emph{Say it:} \emph{maktab}
  \item \emph{Say it:} \emph{maktab}, once more
  \item \emph{Say it:} say \emph{maktab}
  \item \emph{Recall:} say the Arabic for a mosque, then the Arabic for a hotel, then say \emph{maktab} again
\end{itemize}

\begin{tcolorbox}[breakable,colback=teal!4,colframe=teal!35!black,title={Before you move on}]
What does \ar{مكتب} mean? (\textquotedblleft{}An office; a desk\textquotedblright{}.) Say it once more.
\end{tcolorbox}
